\crefalias{section}{appendix}
\section{Algebraic Geometry of the Causal Loop and Complex Multiplication}
\label{app:algebraic_geometry_causal_loop}

This appendix provides the formal mathematical derivation bridging the physical requirement of Unitarity to the algebraic necessity of the Stark-Heegner numbers, as referenced in \cref{sec:fundamental_resonance_filter1}.

\subsection{The Setup: Dimensional Reduction to the Causal Plane}
Information propagation on the $E_8$ substrate ($D_4 \oplus D_4$) is fundamentally local: a signal from one node to another defines a 1-dimensional spatial trajectory evolving over 1-dimensional time, restricting causal propagation to a $(1+1)D$ \textbf{Causal Plane}.

Because the vacuum possesses a fundamental temporal cycle, this sequence of updates must eventually close, forming a periodic loop of length $\Delta$. 

\subsection{Mathematical Translation: The Elliptic Curve}
Topologically, a discrete, periodic $(1+1)D$ causal plane forms a 2-dimensional real torus ($\mathbb{T}^2$). Mathematically, this is equivalent to a 1-dimensional complex torus, known in algebraic geometry as an Elliptic Curve.

The geometry of this causal subtorus is governed by a complex modular parameter $\tau$. As established in Section~\ref{sec:metric_signature}, time acts as the geometric imaginary unit ($i \equiv \omega$) due to the Lorentzian metric signature.

Because the trajectory is physically constrained to the discrete, self-dual structure of the underlying $E_8$ lattice, $\tau$ cannot be a generic continuous complex variable. The self-duality constraint requires the partition function $Z_\Lambda(\tau) = \Theta_\Lambda(\tau)/\eta(\tau)^8$ to be invariant under $S: \tau \to -1/\tau$. The discrete, self-dual lattice structure constrains $\tau$ to a CM point, a fixed point of a M\"obius transformation with integer coefficients acting on the moduli space, requiring $\tau$ to satisfy a quadratic equation $a\tau^2 + b\tau + c = 0$ with $a, b, c \in \mathbb{Z}$ and discriminant $b^2 - 4ac < 0$.

The Unitarity and class number constraints restrict the CM discriminant to the Stark-Heegner set $D \in \{1, 2, 3, 7, 11, 19, 43, 67, 163\}$. The Causality ($\Delta > N = 32$) and Temporal Atomicity ($\Delta \leq 2N = 64$) constraints independently select $\Delta = 43$ as the unique temporal period. 

The identification $D = \Delta$ is a consistency condition: the temporal period must be a valid CM discriminant for the causal loop to admit unitary structure. That $\Delta = 43$ lies within the Stark-Heegner set validates this geometric identification. The self-duality requirement ($\Lambda = \Lambda^*$) equates the lattice with its reciprocal, making the fundamental cycle length and modular period naturally the same integer; the notation $\mathbb{Q}(\sqrt{-\Delta})$ reflects this.

\subsection{The Failure Mode: Isogeny Classes and Causal Degeneracy}
For the system to preserve quantum Unitarity, the history of a state traversing this causal loop must be uniquely reconstructable. Mathematically, this requires the endomorphism ring of the elliptic curve to be a Principal Ideal Domain (PID), which corresponds to a class number of $h=1$.

We analyze the failure mode where $h > 1$. In this regime, the ideal class group of the field is non-trivial. By the theory of Complex Multiplication, this dictates that multiple non-isomorphic elliptic curves will share the exact same endomorphism ring. 

Physically, these isogenous curves represent topologically distinct causal histories. However, because they are locally indistinguishable by their norm structure, a physical measurement at the output node cannot determine which of the $h$ inequivalent curves generated the state. 

\subsection{The Resolution: The Unitarity Constraint}
This local indistinguishability manifests as \textbf{Causal Degeneracy}. The causal history is fundamentally erased, not merely unknown. This ambiguity explicitly violates the unitary requirement that the evolution operator be uniquely invertible ($\hat{U}^\dagger \hat{U} = I$). 

To prevent effective information loss and preserve Unitarity, the universe must strictly forbid non-trivial ideal class groups in its fundamental causal cycles. The algebraic field defining the cycle must possess $h = 1$, restricting the CM discriminant $D$ to the Stark-Heegner numbers. The consistency condition $D = \Delta$ then validates that the temporal period derived from Causality and Atomicity is compatible with unitary causal loop structure.